\documentclass[11pt]{article}

% ----------------------------------------------------------------------------
% Quick selected-method ecology benchmark — experiment report
% Generated by Claude from the on-disk run outputs (seed 42).
% Compile with: pdflatex quick_selected_methods_report.tex
% ----------------------------------------------------------------------------

\usepackage[margin=1in]{geometry}
\usepackage{amsmath}
\usepackage{amssymb}
\usepackage{booktabs}
\usepackage{array}
\usepackage[hidelinks]{hyperref}

\newcommand{\reff}{r_{\mathrm{eff}}}
\newcommand{\Keff}{K_{\mathrm{eff}}}
\newcommand{\rbase}{r_{\mathrm{base}}}

\title{\textbf{A Short Selected-Method Benchmark of General Model-Based\\
Offline RL versus Ecology-Specific Baselines\\
for Adaptive Population Management}}
\author{Experiment report (one seed, demo-scale)}
\date{June 2026}

\begin{document}
\maketitle

\begin{abstract}
We report a short, single-seed benchmark comparing two \emph{general} model-based
offline reinforcement-learning (RL) methods adapted to an ecological
adaptive-management problem---\textbf{MOPO} (pessimistic dynamics ensemble) and
\textbf{RefPlan} (reflect-then-plan)---against two \emph{ecology-specific}
mechanistic baselines---\textbf{PLUS} (BioConserv'18, finite hidden-model Bayesian
planning) and \textbf{MOOR-Ricker} (ExpertSys'23, least-squares mechanistic
offline RL). The study has two settings: a \emph{primary} setting in which the
true environment is a Ricker population model, and a \emph{model-mismatch stress}
setting in which the true environment is a Schaefer/logistic model while PLUS and
MOOR-Ricker retain their Ricker-family assumption. On the primary Ricker task the
two ecology baselines lead on cumulative reward
($16.44$), MOPO is competitive ($16.01$), and RefPlan trails ($14.95$); all
methods keep abundance healthy. In the stress setting the ordering is unchanged
and the intended advantage of the model-agnostic methods does \emph{not} appear:
the Ricker-assuming baselines are barely harmed by the mismatch ($-0.11$), whereas
the uncertainty-aware learned methods lose \emph{more} reward (MOPO $-0.42$,
RefPlan $-0.26$). We argue
this is a property of the chosen dynamics---Ricker and this Schaefer model share
the same equilibrium, so the model family barely affects the optimal policy---and
not a defect of the methods. The only model-agnostic advantage we observe is
\emph{diagnostic}: RefPlan's epistemic uncertainty rises sharply under mismatch.
We also document a compute finding: MOPO's full-budget planning evaluation is
very expensive ($\sim$6.25M single-sample ensemble queries), taking $\sim$2.8--2.9\,h
on an NVIDIA L40S; the Schaefer run first \emph{timed out} at a 4\,h wall on a
slower Tesla~T4 and was completed by re-running the identical configuration on an
L40S. This is a framework/demonstration result, not a tuned scientific benchmark.
\end{abstract}

\section{Introduction and Problem Setting}

We study offline, model-based adaptive management of a single harvested/managed
population. The decision problem is a finite Markov decision process (MDP) with
exact (noise-free) discrete state observation.

\paragraph{State and dynamics.}
The continuous abundance $s_t$ evolves under a population model with a
\emph{hidden} per-episode intrinsic growth rate $\rbase \sim \mathcal{U}(0.95,
1.00)$, drawn at episode reset and held fixed within the episode (never logged).
The agent observes only the discrete abundance bin
$x_t = \min(\lfloor s_t / w \rfloor, N-1)$ with bin width $w = 10$, $N = 100$
bins, and maximum abundance $1000$. Management actions apply additive
perturbations to the growth rate and carrying capacity,
\begin{equation}
  \reff = \rbase + \Delta r(a), \qquad
  \Keff = K_{\mathrm{base}} + \Delta K(a), \qquad K_{\mathrm{base}} = 500,
\end{equation}
and the population then updates under the true model:
\begin{align}
  \text{Ricker (primary):}   \quad & s_{t+1} = s_t \, \exp\!\big(\reff\,(1 - s_t/\Keff)\big), \\
  \text{Schaefer (stress):}  \quad & s_{t+1} = s_t + \reff\, s_t\,(1 - s_t/\Keff),
\end{align}
both clipped to $[0, 1000]$. Episodes are fixed-horizon (length 50); bin~0 is not
absorbing (the population can recover).

\paragraph{Actions.}
Five structured management actions are available in every state
(Table~\ref{tab:actions}); a negative cost denotes revenue.

\begin{table}[h]
\centering
\small
\begin{tabular}{clrrr}
\toprule
id & name & $\Delta r$ & $\Delta K$ & cost \\
\midrule
0 & Do Nothing                 &  0.00 &   0 &  0.00 \\
1 & Aggressive Harvest         & -0.02 &   0 & -0.40 \\
2 & Predator / Disease Control &  0.01 &   0 &  0.20 \\
3 & Intensive Restoration      &  0.00 & 200 &  0.30 \\
4 & Flagship Conservation      &  0.02 & 200 &  0.60 \\
\bottomrule
\end{tabular}
\caption{Action metadata, shared identically by the Ricker and Schaefer
environments.}
\label{tab:actions}
\end{table}

\paragraph{Reward.}
A single first-class reward contract is shared by the environment and all
planners,
\begin{equation}
  R(x_t, a_t) = \alpha \, \frac{x_t}{x_{\max}} - \mathrm{cost}(a_t),
  \qquad \alpha = 1, \; x_{\max} = 100,
\label{eq:reward}
\end{equation}
i.e.\ reward is keyed on the \emph{current} discrete abundance and the action
cost. The objective rewards keeping abundance high while avoiding costly
interventions.

\section{Methods Compared}

We deliberately benchmark only four selected methods, not the full implemented
ladder.

\paragraph{General model-based offline RL (learn the dynamics from data).}
\begin{itemize}
  \item \textbf{MOPO}: a categorical dynamics \emph{ensemble} (5 networks,
  frame-stack 2, bootstrap ratio 0.8, 20 training epochs) combined with a
  \emph{pessimistic} receding-horizon planner. For each candidate action the
  planner simulates the ensemble forward, penalising the discounted return by the
  ensemble disagreement: $\hat r(s,a) = \mathbb{E}[\text{reward}] - \lambda
  \cdot \mathrm{Var}_i(\mathbb{E}_i[x'])$ with $\lambda = 1.0$, planning horizon
  5, and 20 Monte-Carlo rollouts.
  \item \textbf{RefPlan}: a reflect-then-plan tabular model-based planner
  (ensemble size 10, planning horizon 4, 64 candidate sequences, 4 latent
  samples, $\kappa = 5.0$, uncertainty penalty $0.10$). It plans over a learned
  latent posterior and is explicitly uncertainty-aware.
\end{itemize}

\paragraph{Ecology-specific mechanistic baselines (assume the Ricker family).}
\begin{itemize}
  \item \textbf{PLUS} (BioConserv'18): a finite hidden-model Bayesian baseline.
  It maintains a posterior over 21 candidate values of $\rbase$ spanning the
  prior, builds one transition matrix per candidate, solves each by exact value
  iteration ($\gamma = 0.95$), and acts by posterior-averaged Q-values. It
  performs \emph{no} neural training.
  \item \textbf{MOOR-Ricker} (ExpertSys'23): a mechanistic offline RL baseline
  that \emph{learns} the growth parameters $(r, K)$ from the offline transitions
  by least squares (L-BFGS with restarts), discretises the learned Ricker model,
  and plans by value iteration ($\gamma = 0.95$).
\end{itemize}

\section{Experimental Setup}

\subsection{Data, evaluation, and fairness}
A single offline dataset of $25{,}000$ transitions was generated per setting
under a \emph{random} collection policy with seed 42. MOPO, RefPlan and
MOOR-Ricker consume the \emph{same} saved dataset for a given setting; PLUS does
not train from data but is evaluated on the same environment, seeds, episode
count and horizon. All methods are evaluated on 50 episodes of horizon 50 with
discount $0.95$, using an identical sequence of per-episode environment seeds
($42,\dots,91$), so every method sees the same evaluation episodes.

\subsection{The two settings and the mismatch mechanism}
\begin{description}
  \item[Primary (Ricker).] True environment, data and evaluation are all Ricker.
  PLUS uses the true Ricker transition (conditioned on each candidate $\rbase$);
  MOOR-Ricker is well-specified.
  \item[Stress (Schaefer mismatch).] True environment, data and evaluation are
  all Schaefer/logistic, with identical states, actions and reward
  (Eq.~\ref{eq:reward}). MOPO and RefPlan run \emph{unchanged} and learn from the
  Schaefer data. The mechanistic baselines are deliberately kept on their Ricker
  assumption: PLUS sets \verb|use_env_transition=false|, which builds its
  candidate transition matrices from the internal Ricker formula
  $s\,\exp(\reff(1-s/\Keff))$ rather than the true Schaefer transition;
  MOOR-Ricker fits a Ricker model ($\texttt{dynamics\_model=ricker}$) to the
  Schaefer-generated data.
\end{description}
We verified on disk that the Schaefer dataset metadata records
\texttt{active\_env: schaefer}, that the Schaefer PLUS run received
\texttt{model.use\_env\_transition=false}, and that MOPO/RefPlan/MOOR shared the
single Schaefer dataset. The intended hypothesis is that the model-agnostic
methods, which never assume a population-model family, should remain competitive
when the assumed ecological model is wrong.

\subsection{Hardware and scheduling}
Runs were submitted to Slurm as a four-task array (concurrency cap 4). The Ricker
methods executed on NVIDIA L40S nodes; the Schaefer methods executed on a Tesla
T4 node. MOPO and RefPlan used \texttt{device=cuda}; PLUS and MOOR-Ricker used
\texttt{device=cpu} within GPU jobs. A lightweight wrapper recorded wall-clock
runtime, host, GPU name and a GPU-memory sample per method.

\section{Results}

\subsection{Primary setting (Ricker)}
All four methods completed. Table~\ref{tab:ricker} reports the results, sorted by
mean cumulative reward.

\begin{table}[h]
\centering
\small
\begin{tabular}{lcccccc}
\toprule
Method & Cum.\ reward & $[\min,\max]$ & Final & Min & Catastrophic & Runtime \\
       & (mean\,$\pm$\,std) &        & abund. & abund. & rate & (s) \\
\midrule
MOOR-Ricker & $16.436 \pm 0.423$ & $[14.69, 16.87]$ & 49.5 & 38.2 & 0.06 & 158 \\
PLUS        & $16.436 \pm 0.423$ & $[14.69, 16.87]$ & 49.5 & 38.2 & 0.06 & 253 \\
MOPO        & $16.013 \pm 0.680$ & $[13.66, 16.69]$ & 49.5 & 37.8 & 0.06 & 10206 \\
RefPlan     & $14.947 \pm 0.763$ & $[12.61, 16.20]$ & 49.4 & 38.2 & 0.06 & 242 \\
\bottomrule
\end{tabular}
\caption{Primary Ricker setting, 50 evaluation episodes, seed 42. All runs on an
NVIDIA L40S. ``Catastrophic rate'' is the fraction of episodes whose minimum
abundance bin $\le 5$.}
\label{tab:ricker}
\end{table}

\subsection{Stress setting (Schaefer mismatch)}
All four methods are reported (Table~\ref{tab:schaefer}). MOPO required two
attempts: the original array task ran on a slower Tesla~T4 and hit the 4\,h Slurm
wall (state \texttt{TIMEOUT}, elapsed $04{:}01{:}25$) with no CSV---its dynamics
ensemble had trained and checkpointed in $\sim$3\,min, so the timeout occurred
entirely inside the planning evaluation. Re-running the \emph{identical}
full-budget configuration on an NVIDIA~L40S (job 56111492) completed in
$\sim$2.9\,h (\S\ref{sec:compute}). MOPO's Schaefer result is therefore full-budget
and directly comparable to the other cells.

\begin{table}[h]
\centering
\small
\begin{tabular}{lcccccc}
\toprule
Method & Cum.\ reward & $[\min,\max]$ & Final & Min & Catastrophic & Runtime \\
       & (mean\,$\pm$\,std) &        & abund. & abund. & rate & (s) \\
\midrule
MOOR-Ricker & $16.323 \pm 0.480$ & $[14.24, 16.66]$ & 50.0 & 32.7 & 0.08 & 24 \\
PLUS        & $16.323 \pm 0.480$ & $[14.24, 16.66]$ & 50.0 & 32.7 & 0.08 & 127 \\
MOPO        & $15.595 \pm 1.088$ & $[10.99, 16.62]$ & 50.0 & 34.9 & 0.08 & 10365 \\
RefPlan     & $14.688 \pm 0.835$ & $[11.86, 15.90]$ & 49.7 & 32.2 & 0.08 & 324 \\
\bottomrule
\end{tabular}
\caption{Schaefer model-mismatch stress setting, 50 evaluation episodes, seed 42.
PLUS, MOOR-Ricker and RefPlan ran on a Tesla T4; MOPO is the full-budget re-run on
an NVIDIA L40S (the original Tesla-T4 array task timed out at the 4\,h wall;
\S\ref{sec:compute}). PLUS and MOOR-Ricker retain a Ricker-family assumption
against a Schaefer truth.}
\label{tab:schaefer}
\end{table}

\subsection{Epistemic uncertainty and policy entropy}
Table~\ref{tab:unc} reports two auxiliary signals: the mean per-step epistemic
uncertainty each method logs, and the mean policy entropy. These separate the
methods far more than reward does.

\begin{table}[h]
\centering
\small
\begin{tabular}{lcccc}
\toprule
 & \multicolumn{2}{c}{Ricker (primary)} & \multicolumn{2}{c}{Schaefer (stress)} \\
\cmidrule(lr){2-3}\cmidrule(lr){4-5}
Method & uncertainty & entropy & uncertainty & entropy \\
\midrule
MOOR-Ricker & 0.0000 & 0.007 & 0.0000 & 0.014 \\
PLUS        & 0.0028 & 0.015 & 0.0039 & 0.024 \\
MOPO        & 0.0425 & 0.249 & 0.1468 & 0.422 \\
RefPlan     & 0.0668 & 0.370 & 0.2207 & 0.520 \\
\bottomrule
\end{tabular}
\caption{Mean logged epistemic uncertainty and mean policy entropy. Both learned
methods become much less certain under the Schaefer mismatch (RefPlan
$0.067\to0.221$, max $2.25\to6.46$; MOPO $0.043\to0.147$, max $1.56\to5.02$),
while the mechanistic baselines stay confident.}
\label{tab:unc}
\end{table}

\subsection{Compute and runtime}
\label{sec:compute}
Runtime is dominated by MOPO's planning evaluation. The pessimistic planner
performs a genuine receding-horizon rollout in which each candidate action,
rollout and look-ahead step queries the ensemble. The number of single-sample
ensemble \texttt{predict} calls per decision step is
\begin{equation}
  \underbrace{5}_{\text{actions}} \times \underbrace{20}_{\text{rollouts}}
  \times \underbrace{\textstyle\sum_{h=0}^{4}\big(1 + 5\cdot\mathbf{1}[h{<}4]\big)}_{=25}
  = 2500,
\end{equation}
so a 50-episode, 50-step evaluation issues $\approx 6.25\times10^6$ ensemble
forward passes. We measured one \texttt{predict} call at $\sim$28.5\,ms on CPU and
inferred $\sim$1.9\,ms on the L40S (from the observed $\sim$4.0\,s per decision
step). This yields $\sim$2.8\,h for the Ricker MOPO evaluation on the L40S
($\texttt{elapsed}=10206$\,s total, of which $10013.9$\,s was evaluation). The
Schaefer MOPO run first hit a 4\,h \texttt{TIMEOUT} on a slower Tesla~T4, then
completed at full budget on an L40S ($\texttt{elapsed}=10365$\,s; job 56111492).
The remaining methods each finished in under 6\,minutes
(Tables~\ref{tab:ricker},~\ref{tab:schaefer}). GPU-memory accounting was
unreliable for the original batch: the node-level sampler returned $0$\,MB for
every method except a $573$\,MB idle-context reading for Ricker MOPO. A
per-process sampler was subsequently added and \emph{did} work on the Schaefer
MOPO re-run (peak $564$\,MB, process-tree scope), but the other cells predate it.

\section{Analysis and Possible Reasons}

\subsection{Ecology baselines lead; the learned planners pay an approximation tax}
On the primary Ricker task PLUS and MOOR-Ricker achieve the highest reward, MOPO
is close behind, and RefPlan trails. This ordering is consistent with the nature
of the task. The environment is a \emph{well-specified, near-deterministic,
exact-observation finite MDP}: given $\rbase$ the dynamics are deterministic, and
the only stochasticity the agent perceives is the narrow prior over $\rbase$ plus
bin aliasing. PLUS knows the Ricker functional form and the correct
$K_{\mathrm{base}}$, carries a posterior over $\rbase$, and solves each candidate
MDP exactly by value iteration; MOOR-Ricker recovers $(r, K)$ accurately by least
squares and likewise plans by exact value iteration. Both effectively
\emph{solve} the MDP. By contrast, MOPO and RefPlan approximate the dynamics with
a neural ensemble or latent model and plan with short horizons and explicit
pessimism / uncertainty penalties; in a setting where the true model is easy to
identify, that approximation error and built-in conservatism cost reward rather
than protect against it. MOPO outperforming RefPlan is consistent with MOPO's
held-out dynamics accuracy ($0.79$) and its value-based pessimistic rollout
extracting more reward than RefPlan's shorter-horizon, more heavily penalised
($0.10$) latent planning.

\subsection{PLUS and MOOR-Ricker are numerically identical}
In \emph{both} settings PLUS and MOOR-Ricker produce identical per-episode
cumulative rewards (Ricker $16.436$; Schaefer $16.323$), hence identical
abundance statistics. This is not a bug. With exact observations and a shared
reward contract, both methods converge to the same greedy policy: PLUS via a
posterior over $\rbase$ concentrated on the correct Ricker form, MOOR via a
least-squares Ricker fit. They differ only in their auxiliary uncertainty signals
(Table~\ref{tab:unc}), where MOOR's restart disagreement is essentially zero
(well-identified point estimate) and PLUS's posterior spread is small but
non-zero. In this exact-observation regime the two ecology baselines are
effectively redundant; they may diverge under harder (e.g.\ partially observed or
noisier) settings.

\subsection{Why the mismatch stress did not expose the intended advantage}
The central design goal of the stress setting was to reward methods that do
\emph{not} assume a population-model family. Empirically the mismatch did
\emph{not} favour the model-agnostic methods: the Ricker-assuming baselines fell
only $0.11$ ($16.44 \to 16.32$), RefPlan fell $0.26$ ($14.95 \to 14.69$), and MOPO
fell the \emph{most}, $0.42$ ($16.01 \to 15.60$); the within-setting ordering is
unchanged. Two effects combine: (i) the mismatch is largely
\emph{decision-irrelevant} (below), so there is little for a model-agnostic method
to exploit; and (ii) for the uncertainty-aware learned methods it is mildly
\emph{adversarial} on reward (\S\ref{sec:pessimism}). We first attribute the
decision-irrelevance to the choice of dynamics:
\begin{enumerate}
  \item \textbf{Shared equilibrium and local stability.} Both maps have fixed
  points at $s=0$ and $s=K$, and both have the same linearisation at the carrying
  capacity, $f'(K) = 1 - r$, which is stable for $r \in (0,2)$. For
  $\rbase \in [0.95, 1.00]$ both families converge monotonically to the same
  $\Keff$. They differ only in the low-abundance transient (Ricker grows by a
  factor $\approx e^{r}\!\approx\!2.6$ per step, the logistic by
  $\approx 1+r\!\approx\!1.95$).
  \item \textbf{The optimum is family-invariant.} Because reward
  (Eq.~\ref{eq:reward}) rewards sitting near a high abundance and avoiding costly
  actions, and because both families drive the population to the \emph{same}
  $\Keff$ determined by the \emph{same} action metadata, the optimal action
  ranking is essentially identical whether one assumes Ricker or Schaefer. The
  mismatch perturbs the path, not the decision.
  \item \textbf{The baselines are not very ``wrong''.} MOOR re-fits $(r,K)$ to the
  Schaefer data and lands near the true $K$; PLUS uses the correct $K_{\mathrm{base}}$
  and $\rbase$ grid. Only the functional form differs, and the two forms agree at
  the operating point $K$.
  \item \textbf{Coarse discretisation and a narrow prior} (100 bins of width 10;
  $\rbase \in [0.95,1.0]$) further wash out the transient differences between the
  families.
\end{enumerate}
The clearest evidence is that PLUS and MOOR-Ricker remain numerically identical
across both settings: whether the model is the true Ricker env, an internal Ricker
formula, or a Ricker re-fit to Schaefer data, all three collapse to the same
greedy policy---a direct demonstration that the model family does not change the
decision here.

\subsection{Under mismatch, uncertainty rises---aiding calibration but costing reward}
\label{sec:pessimism}
Both learned methods become markedly less certain when the truth becomes Schaefer:
RefPlan's mean epistemic uncertainty more than triples ($0.067 \to 0.221$; max
$2.25 \to 6.46$) and MOPO's more than triples as well ($0.043 \to 0.147$; max
$1.56 \to 5.02$), while the mechanistic baselines stay confident (MOOR $\approx 0$,
PLUS $\approx 0.004$; Table~\ref{tab:unc}). This has two opposite consequences. On
the positive side it is a genuine \emph{calibration} signal: the model-agnostic
methods \emph{detect} that they are operating outside their training distribution,
whereas PLUS and MOOR---using a wrong-but-adequate Ricker model---remain
overconfident. On the negative side it directly costs reward: MOPO's pessimistic
planner subtracts $\lambda\cdot\mathrm{Var}_i$ from every imagined step
($\lambda=1$) and RefPlan applies an explicit uncertainty penalty, so greater
disagreement makes both plan more conservatively. This explains why MOPO loses the
most reward under the mismatch ($-0.42$): the unfamiliar Schaefer dynamics inflate
ensemble disagreement, and pessimism converts that disagreement into caution even
though no catastrophe is actually present in this benign regime. In short, the
mismatch here does not reward model-agnosticism---it \emph{penalises} the
uncertainty-aware methods on reward while improving their calibration.

\subsection{Safety profile}
All methods keep the population healthy in both settings: mean abundance sits near
bin $49$--$50$ (the population equilibrates near $K=500 \Rightarrow$ bin $50$),
and catastrophic episodes are rare ($0.06$ Ricker, $0.08$ Schaefer). The slightly
higher catastrophic rate and lower minimum abundance ($\sim$32 vs.\ $\sim$38)
under Schaefer reflect the logistic's slower recovery from low abundance, but no
method collapses. The task therefore lives in a ``safe'' region where the
abundance-maximising and safety objectives are aligned---another reason the
mismatch is not punishing.

\section{Limitations}
\begin{itemize}
  \item \textbf{Single seed, demo-scale compute, no tuning.} Hyperparameters are
  the handoff defaults. No claim of final superiority is warranted.
  \item \textbf{Schaefer MOPO required a re-run.} The original Tesla-T4 array task
  timed out at the 4\,h wall; the cell was filled by re-running the identical
  full-budget configuration on an L40S ($\sim$2.9\,h). The MOPO result is therefore
  full-budget and comparable, but the two Schaefer-setting runs used different GPUs.
  \item \textbf{Weak stress design.} As analysed above, the Ricker$\leftrightarrow$Schaefer
  swap does not make the model family decision-relevant, so it is not a strong
  test of the intended hypothesis.
  \item \textbf{GPU-memory metric unavailable} for these runs.
  \item \textbf{Reproducibility wrinkle.} The Ricker array ran a slightly earlier
  version of the harness without explicit \texttt{active\_env} /
  \texttt{use\_env\_transition} overrides; the defaults coincided with the correct
  Ricker behaviour, so the results are valid, but the two settings were not
  produced by byte-identical harness code.
\end{itemize}

\section{Designing Scenarios that Favour RefPlan over Fixed-Family Baselines}
\label{sec:design}

RefPlan learns a flexible transition model from data---here a tabular, per-$(x,a)$
ensemble with an epistemic-uncertainty penalty---and commits to \emph{no}
population-model family. PLUS and MOOR-Ricker instead assume the dynamics belong to
one known parametric family (or, more generally, a small known set of families such
as $\{$Ricker, logistic$\}$) and merely fit that family's free parameters. RefPlan
can only out-score such a baseline when the committed family assumption is both
wrong and costly. This section formalises when that happens and gives concrete,
repo-compatible designs.

\subsection{Four conditions for a RefPlan advantage}
A scenario rewards the family-agnostic learner over a baseline that assumes one or
two known families when \emph{all four} of the following hold:
\begin{enumerate}
  \item \textbf{Out-of-family truth.} The true transition law lies outside every
  assumed family. If the baseline assumes $\{$Ricker, logistic$\}$, the truth must
  be a \emph{third} form (or a mixture / regime switch), because any model inside
  the assumed set can simply be fitted. Our Schaefer stress partly failed here:
  the logistic is itself a family one would naturally assume.
  \item \textbf{Decision-relevance.} The misspecification must change the
  \emph{optimal action}, not merely the transient path. Misspecification at the
  operating point is harmless if both models share the equilibrium and the reward
  depends only on it---the central reason the Ricker$\leftrightarrow$Schaefer swap
  was inert (\S\ref{sec:pessimism}).
  \item \textbf{Learnability from data.} RefPlan's tabular model is accurate only
  where the offline data covers the relevant $(x,a)$ cells. The behaviour policy
  must therefore \emph{visit} the decision-critical region (e.g.\ low abundance),
  or the flexible model has nothing to learn there and falls back to its prior.
  \item \textbf{An asymmetric / catastrophic cost.} The advantage is largest when
  getting that region wrong is expensive (collapse, extinction) and the region is
  partly data-sparse, so RefPlan's uncertainty penalty ($\kappa$, the $0.10$
  penalty) steers it away while a confident parametric extrapolation walks in.
\end{enumerate}
Conditions 1--2 make the family assumption \emph{wrong}; conditions 3--4 let the
flexible learner \emph{exploit} it.

\subsection{True-environment families outside Ricker/logistic}
Each keeps the same states, actions, reward and hidden $\rbase$; only the update
map changes (a one-method subclass of \texttt{RickerEnv}, exactly as
\texttt{SchaeferEnv} already is). Choose the true family so that even a
\emph{two}-model $\{$Ricker, logistic$\}$ baseline remains misspecified.
\begin{description}
  \item[Critical depensation / Allee.]
  $s_{t+1} = s_t \exp\!\big(\reff\,(1 - s_t/\Keff)\,(s_t/C - 1)\big)$ with a
  critical threshold $C$: growth is \emph{negative} below $C$, so the population
  collapses once pushed under it. Compensatory families (Ricker, logistic) have no
  such threshold, so a family-assuming agent harvests through $C$ expecting
  recovery. This is the canonical ``wrong model $\Rightarrow$ catastrophic
  over-harvest'' case.
  \item[$\theta$-logistic.]
  $s_{t+1} = s_t + \reff\, s_t\big(1 - (s_t/\Keff)^{\theta}\big)$ with
  $\theta \neq 1$. For $\theta \gg 1$ density dependence concentrates near $K$
  (flat at low density); for $\theta \ll 1$, the opposite. Neither Ricker nor
  logistic ($\theta{=}1$) matches the curvature, and the optimal harvest level
  shifts with $\theta$.
  \item[Regime-switching / mixture.] Each episode (or each step, via a hidden
  environmental state) draws the dynamics from a family \emph{not} in the candidate
  set, or alternates between two such families. A finite-candidate Bayesian
  baseline (PLUS) cannot place posterior mass on a regime it does not enumerate;
  RefPlan learns the marginal transition.
  \item[Delayed density dependence (non-Markovian).] Let growth depend on $s_t$
  \emph{and} $s_{t-1}$ (a lag-1 effect). A first-order Ricker/logistic family is
  structurally blind to the lag, whereas RefPlan/MOPO with frame-stacking observe
  the history and can capture it---decision-relevant when the lag changes optimal
  harvest timing.
\end{description}

\subsection{Experimental knobs that widen the gap}
\begin{itemize}
  \item \textbf{Coverage of the danger zone.} Replace the pure \texttt{random}
  collection policy with one that also visits low abundance (e.g.\ a mix of random
  and a harvesting heuristic), so the collapse/threshold region is in-distribution
  for RefPlan; otherwise condition~3 fails.
  \item \textbf{Sharper reward on collapse.} The current contract
  (Eq.~\ref{eq:reward}) barely penalises low abundance (catastrophic rate
  $6$--$8\%$). Add an explicit penalty or make extinction absorbing below a
  threshold, so a single mismanaged collapse dominates the return---precisely where
  avoiding the out-of-family failure pays.
  \item \textbf{Push into the regime where transient $\neq$ equilibrium.} Widen the
  hidden prior and/or raise $\reff$ toward the Ricker over-compensation/oscillation
  range ($r \gtrsim 2$), so that mispredicting the transient (which the wrong family
  does) changes the chosen action.
  \item \textbf{Process / observation noise.} Heteroscedastic or heavy-tailed noise
  hurts point-estimate baselines (MOOR) and a fixed deterministic candidate grid
  (PLUS) more than an ensemble with an uncertainty penalty; concentrate it where
  mistakes are costly.
  \item \textbf{Longer evaluation horizon}, so that getting the threshold/transient
  right compounds over more steps.
\end{itemize}

\subsection{A concrete protocol}
A minimal high-leverage instance: true environment $=$ \emph{critical-depensation}
Ricker with threshold $C \approx 0.3\,K$; baselines $=$ PLUS and MOOR-Ricker kept on
the compensatory Ricker family (optionally give PLUS a second candidate family,
logistic, to test the ``two known models'' case---both still lack depensation);
MOPO and RefPlan unchanged. Use a collection policy that drives the population near
and below $C$ so the offline data contains collapses, and add an extinction
penalty. Hypothesis: the family-assuming baselines, lacking any representation of
the threshold, over-harvest into collapse on a non-trivial fraction of episodes,
whereas RefPlan---having learned the negative-growth region and penalising
uncertain, data-sparse states---keeps abundance above $C$ and attains higher return
at a lower catastrophic rate.

\subsection{Pitfalls and fairness}
Keep states, actions, reward, seeds, episode count and horizon identical across
methods, and give the parametric baselines their \emph{best} fit---the win must come
from family misspecification, not from under-fitting them. Ensure the deviation is
learnable: if the data never visits the out-of-family region, RefPlan cannot learn
it and the comparison reverts to the inert case. Avoid the opposite failure too:
dynamics so noisy or so data-starved that \emph{every} method fails. As throughout,
firm conclusions require multiple seeds and tuned hyperparameters; a single demo
seed only indicates a direction.

\section{Conclusions and Recommendations}

The framework runs end-to-end and now compares adapted general model-based offline
RL methods (MOPO, RefPlan) against ecology-specific baselines (PLUS, MOOR-Ricker)
under shared reward and dynamics semantics. On reward, the ecology baselines lead
and MOPO is competitive; RefPlan trails. The intended ``model-agnostic methods win
under mismatch'' story is \emph{not} supported by this run: the chosen
Ricker/Schaefer mismatch does not change the optimal policy, and---because both
learned methods are uncertainty-aware---the mismatch actually \emph{penalises} them
on reward (MOPO loses the most, $-0.42$) while improving their calibration (both
MOPO and RefPlan detect the shift via sharply higher uncertainty).

To actually test the hypothesis, the stress setting must make the misspecification
\emph{decision-relevant}, ideally toward catastrophe---the regime the BioConserv'18
/ ExpertSys'23 lineage cares about. Section~\ref{sec:design} gives a detailed recipe
(out-of-family true dynamics such as critical depensation, coverage of the danger
zone, and a collapse-sensitive reward) for constructing scenarios in which RefPlan
should beat baselines that assume one or two known population families.

For compute, MOPO's full-budget planning evaluation should either be given a
longer wall on a fast GPU, run at a (disclosed) reduced planning budget
($\texttt{num\_rollouts}{=}1$, $\texttt{horizon}{=}3$ reduces the cost
$\sim$40$\times$ to $\sim$1.5\,s/step), or---for a real benchmark---accelerated by
vectorising the \texttt{predict} call across ensemble members and rollouts, which
would turn $\sim$28\,ms/call into sub-millisecond and make full budget finish in
minutes. Finally, the study should be repeated over multiple seeds with tuned
hyperparameters before any scientific claim is drawn.

\appendix
\section{Exact configuration}
\small
\noindent\textbf{Environment.} $N=100$ bins, $w=10$, $\max=1000$;
$\rbase\sim\mathcal{U}(0.95,1.00)$; $K_{\mathrm{base}}=500$; reward $\alpha=1$,
$x_{\max}=100$; episode length 50; fixed horizon (no extinction termination).\\
\textbf{Dataset.} random policy, $25{,}000$ transitions, seed 42; shared by
MOPO/RefPlan/MOOR per setting.\\
\textbf{Evaluation.} 50 episodes, horizon 50, $\gamma=0.95$, episode seeds
$42\dots91$ (identical across methods).\\
\textbf{MOPO.} ensemble size 5, frame-stack 2, bootstrap 0.8, 20 epochs;
pessimistic planner $\lambda=1.0$, horizon 5, 20 rollouts.\\
\textbf{RefPlan.} ensemble size 10, planning horizon 4, 64 sequences, 4 latent
samples, $\kappa=5.0$, uncertainty penalty $0.10$, 5 actions.\\
\textbf{PLUS.} 21 candidate models, $\gamma=0.95$, transition grid 2000,
likelihood floor $10^{-12}$; \texttt{use\_env\_transition} $=$ true (Ricker) /
false (Schaefer).\\
\textbf{MOOR-Ricker.} \texttt{dynamics\_model=ricker}, learn $\{r,K\}$, 10
restarts, L-BFGS max-iter 10, transition grid 1000, $\gamma=0.95$.\\
\textbf{Hardware.} Ricker methods: NVIDIA L40S; Schaefer methods: Tesla T4. MOPO,
RefPlan on \texttt{cuda}; PLUS, MOOR on \texttt{cpu}.

\end{document}
